% GENERATED FILE. Edit canonical lessons, then run npm run generate:books.
% canonical-source-hash: fnv1a64:9f4851a9a461d8f7
% canonical-lessons: FA-C357-khamir, FA-C357-varz-dadan, FA-C357-var-amadan, FA-C357-khamirmaye, FA-C357-tanur

\chapter{Dough and bread}
\label{ch:fa-dough-and-bread}
\hlchaptermodality{357}

\begin{chapteropening}
\textbf{By the end of this chapter:} \emph{I can talk about dough, kneading it, letting it rise, yeast and the clay oven.}

Complete the last lesson of chapter 357: I can talk about dough, kneading it, letting it rise, yeast and the clay oven.
\end{chapteropening}

\section[\texorpdfstring{\fa{خمیر} (khamir)}{خمیر (khamir)}]{\fa{خمیر} (khamir) — dough}

\label{lesson:FA-C357-khamir}



\begin{quote}
Before the new one: say the Persian for to sauté, then the Persian for fried onion.
\end{quote}

\subsection*{You'll want to know: \fa{خمیر}}
\textbf{\fa{خمیر}} (\emph{khamir}) — \textquotedblleft{}dough\textquotedblright{}.

\begin{quote}
\fa{نانوا} \fa{خمیر} \fa{را} \fa{روی} \fa{میز} \fa{پهن} \fa{کرد}.
\end{quote}

\begin{quote}
\emph{Nânvâ khamir râ ru-ye miz pahn kard.}
\end{quote}

\begin{quote}
The baker spread the dough out on a table.
\end{quote}

Also any paste. \emph{Khamir-e nân} \textquotedblleft{}bread dough\textquotedblright{}.

\begin{cousinweb}
\emph{khamir} is borrowed from Arabic \emph{khamīr}, \textquotedblleft{}leavened dough\textquotedblright{} (EtymDB).
\end{cousinweb}

\begin{grammarlens}[title={What you've built}]
One more word for telling, explaining and describing at B1.
\end{grammarlens}

\subsection*{Your turn}
\begin{itemize}
\raggedright
  \item \emph{Say it:} \emph{khamir}
  \item \emph{Say it:} \emph{khamir}, once more
  \item \emph{Say it:} a sentence of your own with \emph{khamir}
  \item \emph{Recall:} say the Persian for to sauté, then the Persian for fried onion, then say \emph{khamir} again
\end{itemize}

\begin{tcolorbox}[breakable,colback=teal!4,colframe=teal!35!black,title={Before you move on}]
What does \fa{خمیر} mean? (\textquotedblleft{}Dough\textquotedblright{}.) Say it once more.
\end{tcolorbox}
\section[\texorpdfstring{\fa{ورز} \fa{دادن} (varz dâdan)}{ورز دادن (varz dâdan)}]{\fa{ورز} \fa{دادن} (varz dâdan) — to knead}

\label{lesson:FA-C357-varz-dadan}



\begin{quote}
Before the new one: say the Persian for fried onion, then the Persian for dough.
\end{quote}

\subsection*{You'll want to know: \fa{ورز} \fa{دادن}}
\textbf{\fa{ورز} \fa{دادن}} (\emph{varz dâdan}) — \textquotedblleft{}to knead\textquotedblright{}.

\begin{quote}
\fa{او} \fa{خمیر} \fa{را} \fa{ده} \fa{دقیقه} \fa{ورز} \fa{داد}.
\end{quote}

\begin{quote}
\emph{U khamir râ dah daqiqe varz dâd.}
\end{quote}

\begin{quote}
She kneaded the dough for ten minutes.
\end{quote}

Present \emph{varz midahad}.

\begin{cousinweb}
\emph{varz} \textquotedblleft{}working, toil\textquotedblright{} is a native Persian word, also seen in \emph{varzesh} \textquotedblleft{}sport\textquotedblright{}; \emph{dâdan} \textquotedblleft{}to give\textquotedblright{} is inherited Persian. EtymDB lists no entry.
\end{cousinweb}

\begin{grammarlens}[title={What you've built}]
One more word for telling, explaining and describing at B1.
\end{grammarlens}

\subsection*{Your turn}
\begin{itemize}
\raggedright
  \item \emph{Say it:} \emph{varz dâdan}
  \item \emph{Say it:} \emph{varz dâdan}, once more
  \item \emph{Say it:} a sentence of your own with \emph{varz dâdan}
  \item \emph{Recall:} say the Persian for fried onion, then the Persian for dough, then say \emph{varz dâdan} again
\end{itemize}

\begin{tcolorbox}[breakable,colback=teal!4,colframe=teal!35!black,title={Before you move on}]
What does \fa{ورز} \fa{دادن} mean? (\textquotedblleft{}To knead\textquotedblright{}.) Say it once more.
\end{tcolorbox}
\section[\texorpdfstring{\fa{ور} \fa{آمدن} (var âmadan)}{ور آمدن (var âmadan)}]{\fa{ور} \fa{آمدن} (var âmadan) — to rise (of dough)}

\label{lesson:FA-C357-var-amadan}



\begin{quote}
Before the new one: say the Persian for dough, then the Persian for to knead.
\end{quote}

\subsection*{You'll want to know: \fa{ور} \fa{آمدن}}
\textbf{\fa{ور} \fa{آمدن}} (\emph{var âmadan}) — \textquotedblleft{}to rise (of dough)\textquotedblright{}.

\begin{quote}
\fa{خمیر} \fa{در} \fa{جای} \fa{گرم} \fa{ور} \fa{آمد}.
\end{quote}

\begin{quote}
\emph{Khamir dar jâ-ye garm var âmad.}
\end{quote}

\begin{quote}
The dough rose in a warm place.
\end{quote}

Present \emph{var miâyad}. \emph{Var} is an old word for \textquotedblleft{}up\textquotedblright{}.

\begin{cousinweb}
\emph{var} \textquotedblleft{}up, on\textquotedblright{} is an old native Persian word, and \emph{âmadan} \textquotedblleft{}to come\textquotedblright{} is inherited Persian: to come up.
\end{cousinweb}

\begin{grammarlens}[title={What you've built}]
One more word for telling, explaining and describing at B1.
\end{grammarlens}

\subsection*{Your turn}
\begin{itemize}
\raggedright
  \item \emph{Say it:} \emph{var âmadan}
  \item \emph{Say it:} \emph{var âmadan}, once more
  \item \emph{Say it:} a sentence of your own with \emph{var âmadan}
  \item \emph{Recall:} say the Persian for dough, then the Persian for to knead, then say \emph{var âmadan} again
\end{itemize}

\begin{tcolorbox}[breakable,colback=teal!4,colframe=teal!35!black,title={Before you move on}]
What does \fa{ور} \fa{آمدن} mean? (\textquotedblleft{}To rise (of dough)\textquotedblright{}.) Say it once more.
\end{tcolorbox}
\section[\texorpdfstring{\fa{خمیرمایه} (khamirmâye)}{خمیرمایه (khamirmâye)}]{\fa{خمیرمایه} (khamirmâye) — yeast, leaven}

\label{lesson:FA-C357-khamirmaye}



\begin{quote}
Before the new one: say the Persian for to knead, then the Persian for to rise.
\end{quote}

\subsection*{You'll want to know: \fa{خمیرمایه}}
\textbf{\fa{خمیرمایه}} (\emph{khamirmâye}) — \textquotedblleft{}yeast, leaven\textquotedblright{}.

\begin{quote}
\fa{او} \fa{کمی} \fa{خمیرمایه} \fa{در} \fa{آب} \fa{گرم} \fa{حل} \fa{کرد}.
\end{quote}

\begin{quote}
\emph{U kami khamirmâye dar âb-e garm hall kard.}
\end{quote}

\begin{quote}
She dissolved a little yeast in warm water.
\end{quote}

\emph{khamir} \textquotedblleft{}dough\textquotedblright{} + \emph{mâye} \textquotedblleft{}substance, starter\textquotedblright{}.

\begin{cousinweb}
\emph{khamirmâye} joins \emph{khamir} \textquotedblleft{}dough\textquotedblright{}, borrowed from Arabic \emph{khamīr} (EtymDB), and the native Persian \emph{mâye} \textquotedblleft{}substance, starter\textquotedblright{}.
\end{cousinweb}

\begin{grammarlens}[title={What you've built}]
One more word for telling, explaining and describing at B1.
\end{grammarlens}

\subsection*{Your turn}
\begin{itemize}
\raggedright
  \item \emph{Say it:} \emph{khamirmâye}
  \item \emph{Say it:} \emph{khamirmâye}, once more
  \item \emph{Say it:} a sentence of your own with \emph{khamirmâye}
  \item \emph{Recall:} say the Persian for to knead, then the Persian for to rise, then say \emph{khamirmâye} again
\end{itemize}

\begin{tcolorbox}[breakable,colback=teal!4,colframe=teal!35!black,title={Before you move on}]
What does \fa{خمیرمایه} mean? (\textquotedblleft{}Yeast, leaven\textquotedblright{}.) Say it once more.
\end{tcolorbox}
\section[\texorpdfstring{\fa{تنور} (tanur)}{تنور (tanur)}]{\fa{تنور} (tanur) — a clay oven, a bread oven}

\label{lesson:FA-C357-tanur}



\begin{quote}
Before the new one: say the Persian for to rise, then the Persian for yeast.
\end{quote}

\subsection*{You'll want to know: \fa{تنور}}
\textbf{\fa{تنور}} (\emph{tanur}) — \textquotedblleft{}a clay oven, a bread oven\textquotedblright{}.

\begin{quote}
\fa{نان} \fa{تازه} \fa{را} \fa{از} \fa{تنور} \fa{بیرون} \fa{آورد}.
\end{quote}

\begin{quote}
\emph{Nân-e tâze râ az tanur birun âvard.}
\end{quote}

\begin{quote}
He took the fresh bread out of the clay oven.
\end{quote}

A round oven of clay or sunk in the ground; \emph{nân-e tanuri} \textquotedblleft{}oven bread\textquotedblright{}.

\begin{cousinweb}
\emph{tanur} is inherited from Middle Persian \emph{tanūr}, \textquotedblleft{}oven\textquotedblright{} (EtymDB), an old word that EtymDB takes back to Akkadian; Arabic and Aramaic have it too.
\end{cousinweb}

\begin{grammarlens}[title={What you've built}]
One more word for telling, explaining and describing at B1.
\end{grammarlens}

\subsection*{Your turn}
\begin{itemize}
\raggedright
  \item \emph{Say it:} \emph{tanur}
  \item \emph{Say it:} \emph{tanur}, once more
  \item \emph{Say it:} a sentence of your own with \emph{tanur}
  \item \emph{Recall:} say the Persian for to rise, then the Persian for yeast, then say \emph{tanur} again
\end{itemize}

\begin{tcolorbox}[breakable,colback=teal!4,colframe=teal!35!black,title={Before you move on}]
What does \fa{تنور} mean? (\textquotedblleft{}A clay oven, a bread oven\textquotedblright{}.) Say it once more.
\end{tcolorbox}
